\documentclass[10pt,openbib,a4paper,openright,a4paper]{article}%
\usepackage{amsfonts}
\usepackage[utf8]{inputenc}
\usepackage{graphicx}
\usepackage{eurosym}
\usepackage{color}
\usepackage{xcolor}
\usepackage[french]{babel}
\usepackage[french]{minitoc}
\usepackage{fancyhdr,amsmath,amsthm,amssymb,fancybox}
\usepackage{calc}
\usepackage[a4paper]{geometry}
\usepackage{mathabx}
%\usepackage[dvipsnames]{xcolor}
\textwidth 16.5cm \oddsidemargin 0.0cm \evensidemargin 0.2cm
\topmargin -2.5cm \textheight 26cm

%
%\hfill \rule{0.3cm}{0.3cm}
%-----------------------------------------
\vfuzz2pt \hfuzz2pt
% THEOREMS ------------------------------
\newtheorem{thf}{{Théorème}}
\newtheorem{corf}{{Corollaire}}
\newtheorem{propf}{\underline{Proposition}}
\newtheorem{prodef}{\underline{Proposition-définition}}
\newtheorem{lemf}{\underline{Lemme}}
\newtheorem{thm}{\underline{Théorème}}
\newtheorem{cor}{\underline{Corollaire}}
\newtheorem{lem}{\underline{Lemme}}
%\newtheorem{prop}{\underline{Proposition}}
\newtheorem{defn}{\underline{Définition}}
\newtheorem{deff}{\underline{Définitions}}
\newtheorem{remfs}{\underline{Remarques}}
\newtheorem{remf}{\underline{Remarque}}
\newtheorem{exercice}{Exercice}
\newtheorem{exemple}{Exemple}
\newtheorem{qcm}{{\underline{Q.C.M}}}
\newtheorem{prob}{{\underline{Problème}}}
\newtheorem{propri}{{\underline{Propriétés}}}
\newtheorem{appli}{\underline{Application}}

%\newtheorem{preuve}[Proof]{Preuve}
%\newenvironnement{preuve}[1][Proof]{Preuve}{\hfill\ \rule{0.5em}{0.5em}}
% MATH  -----------------------------------
\newcommand{\norm}[1]{\left\Vert#1\right\Vert}
%\newcommand{\abs}[1]{\left\vert#1\right\vert}
\newcommand{\set}[1]{\left\{#1\right\}}
\newcommand{\ensR}{\mathbb R}
\newcommand{\ensN}{\mathbb N}
\newcommand{\ensQ}{\mathbb Q}
\newcommand{\ensZ}{\mathbb Z}
\newcommand{\ensC}{\mathbb C}
%\newcommand{\eps}{\varepsilon}
\newcommand{\To}{\longrightarrow}
\newcommand{\orange}[1]{\textcolor{orange}{#1}}
\newcommand{\borange}[1]{\textbf{\textcolor{orange}{#1}}}
\newcommand{\BX}{\mathbf{B}(X)}
\newcommand{\A}{\mathcal{A}}
\newcommand{\biindice}[3]
{
	\renewcommand{\arraystretch}{0.7}
	\begin{array}[t]{c}
		{\displaystyle #1}\\ {\scriptstyle #2}\\ {\scriptstyle #3}
	\end{array}
	\renewcommand{\arraystretch}{1}
}
\newcommand{\triindice}[4]
{
	\renewcommand{\arraystretch}{0.7}
	\begin{array}[t]{c}
		{\displaystyle #1}\\ {\scriptstyle #2}\\ {\scriptstyle
			#3}\\{\scriptstyle #4}
	\end{array}
	\renewcommand{\arraystretch}{1}
}
\def\pointir{\unskip . ---  \ignorespaces}
\def\tir{\unskip --\ignorespaces}
%\newtheorem{preuve}{Preuve}%{\hfill \rule{0.5cm}{0.5cm}}
% -------------------------------------------
\colorlet{bordeaux}{red!60!blue!20!darkgray}

\begin{document}
	%\tableofcontents
	%\dominitoc
	%\chapter{Fonction à une seule variable réelle}
	%\minitoc
	
	
	
	\newenvironment{correct}%
	{\noindent\ignorespaces}%
	{\par\noindent\ignorespacesafterend}
	
\includegraphics{algebre25.png}
	
	\textbf{{\large Exercice 1 : Applications} (CLO 1, 2) 4 points}
	\begin{enumerate}
		
		\item Les applications suivantes sont-elles injectives ? sont-elles surjectives ? \hfill 1.5pt
		\begin{enumerate}
			\item $f : \ensR \longrightarrow \ensR$ , $f(x) = x^2$. 
			\item $g : \ensR \longrightarrow \ensR$ , $f(x) = x^3$.
			\item $h : \ensR \longrightarrow \ensZ$ , $f(x) = E(x)$ (où $E$ représente la partie entière d'un réel). 
		\end{enumerate}
		\item Soient $f$ et $g$ deux applications de $\ensR \longrightarrow \ensR$, montrer que : 
		\begin{enumerate}
			\item $g \circ f$ est injective $\mapsto$ $f$ est injective. \hfill 1.25pt
			\item $g \circ f$ est surjective $\mapsto$ $g$ est surjective. \hfill 1.25pt
		\end{enumerate} 

	
	\end{enumerate}
	
	\textbf{{\large Exercice 2 : Relation et classe d'équivalence} (CLO3) 3.5 points}
	
	\begin{enumerate}
		\item Rappeler les caractéristiques d'une relation d'équivalence. \hfill 0.5pt
		\item Soit $\mathcal{R}$ la relation définie sur $\ensR$ par : 
		$$ x \mathcal{R} y \Leftrightarrow f(x) = f(y)$$ où $f(t) = \frac{t^3 + 2}{t^2 + 1}$. 
		\begin{enumerate}
			\item Montrer que $\mathcal{R}$ est une relation d'équivalence. \hfill 1.5pt
			\item Donner le tableau de variation de la fonction $f$. \hfill 0.5pt 
			\item En déduire à partir du tableau de variation de $f$ le cardinal de $cl(a)$ la classe d'équivalence d'un réel $a$. \hfill 1 pt
		\end{enumerate}
	\end{enumerate}
	
\textbf{{\large Exercice 3 : Groupes symétriques} (CLO4, 5) 5 points}

\medskip
\textbf{A. Questions à choix multiples (QCM) 2points}

\begin{enumerate}
\item Soit $\sigma \in S_n$. Une \textbf{inversion} d’une permutation $\sigma$ est :
\begin{enumerate}
\item un couple $(i,j)$ tel que $i<j$ et $\sigma(i)<\sigma(j)$ ;
\item un couple $(i,j)$ tel que $i<j$ et $\sigma(i)>\sigma(j)$ ;
\item un élément $i$ tel que $\sigma(i)=i$ ;
\item une transposition apparaissant dans la décomposition de $\sigma$.
\end{enumerate}

\item Un \textbf{cycle de longueur $p$} est :
\begin{enumerate}
\item une permutation ayant exactement $p$ points fixes ;
\item une permutation qui est le produit de $p$ transpositions ;
\item une permutation qui permute cycliquement $p$ éléments distincts et fixe les autres ;
\item une permutation paire.
\end{enumerate}

\item Une \textbf{transposition} est :
\begin{enumerate}
\item une permutation qui échange deux éléments et fixe tous les autres ;
\item une permutation de signature $-1$ ;
\item un cycle de longueur $3$ ;
\item une permutation possédant exactement une inversion.
\end{enumerate}

\item La \textbf{signature} d’une permutation $\sigma$ est :
\begin{enumerate}
\item le nombre total de cycles dans la décomposition de $\sigma$ ;
\item égale à $+1$ si $\sigma$ est paire et à $-1$ si $\sigma$ est impaire ;
\item le nombre d’inversions de $\sigma$ ;
\item toujours égale à $1$.
\end{enumerate}
\end{enumerate}

\textbf{B. Soient les permutations suivantes :}

\[
\tau_1 =
\begin{pmatrix}
1 & 2 & 3 & 4 & 5 & 6 & 7 \\
2 & 3 & 1 & 5 & 7 & 4 & 6
\end{pmatrix}
\]

\begin{enumerate}
\item Déterminer le nombre d’inversions et la parité de $\tau_1$. \hfill 1pt

\item Décomposer la permutations $\tau_1$ en produits de cycles à supports disjoints. \hfill 0.5pt

\item Décomposer la permutation $\tau_1$ en produits de transpositions. \hfill 0.5pt

\item Calculer $\tau_1^{100}$. \hfill 1pt
\end{enumerate}


\textbf{{\large Exercice 4 : Groupes et anneaux} (CLO4,6) 7.5 points}

\textit{Les questions sont indépendantes.} 
\begin{enumerate}
	\item \textbf{QCM 2 points}

\begin{enumerate}

\item Un \textbf{groupe cyclique} est :
\begin{enumerate}
\item un groupe commutatif fini ;
\item un groupe engendré par un seul élément ;
\item un groupe dont tous les éléments sont d’ordre fini ;
\item un groupe ayant exactement un générateur.
\end{enumerate}

\item Soit $(G,\cdot)$ un groupe et $a\in G$. Le sous-ensemble
$\langle a\rangle = \{a^n \mid n\in\mathbb{Z}\}$ est :
\begin{enumerate}
\item toujours un sous-groupe de $G$ ;
\item un sous-groupe si et seulement si $G$ est abélien ;
\item égal à $G$ pour tout $a\in G$ ;
\item un anneau commutatif.
\end{enumerate}

\item Quelle affirmation est correcte ?
\begin{enumerate}
\item Tout groupe monogène est cyclique ;
\item Tout groupe cyclique est non abélien ;
\item Tout groupe monogène est fini ;
\item Aucun groupe cyclique n’est abélien.
\end{enumerate}

\item Un \textbf{anneau} $(A,+,\cdot)$ est une structure telle que :
\begin{enumerate}
\item $(A,+)$ est un groupe abélien et la multiplication est associative ;
\item $(A,+)$ et $(A,\cdot)$ sont deux groupes ;
\item la multiplication est toujours commutative ;
\item tout élément admet un inverse multiplicatif.
\end{enumerate}

\end{enumerate}
	\item \textbf{Racine troisième}
	\begin{enumerate}
		\item Montrer que $ U_3 = \{ z  \in \ensC / z^3 = 1\}$ est un sous-groupe de $(\ensC^*, \times)$. \hfill 1pt
		\item Montrer que $U_3$ est un groupe monogène. Est-t-il cyclique ? \hfill 1pt
	\end{enumerate} 
	
	\item \textbf{Morphisme de groupes}
	\\Soit $(G,*)$ et $(G', T)$ deux groupes. Soient H et K deux sous-groupes de G et G' respectivement. \\ Soit $f : G \longrightarrow G'$ un morphisme de groupes. 
	\begin{enumerate}
		\item Montrer que $f(H)$ est un sous groupe de $G'$. \hfill 1pt 
		\item Montrer que $f^{-1}(K)$ est un sous groupe de $G$. \hfill 1pt
	\end{enumerate}
	\item Soit $(\ensZ^2, + , *)$ tel que : 
	$$ (a,b)+(c,d) = (a+c, b+d)$$
	$$ (a,b)*(c,d) = (ac, ad+bc)$$
 Montrer que $(\ensZ^2, + , *)$ est un anneau commutatif. \hfill 1.5pt
		
\end{enumerate}

\end{document}